\documentclass[11pt]{article}
% Preambulo comun para los informes de la ronda 3.
% Se incluye con % Preambulo comun para los informes de la ronda 3.
% Se incluye con % Preambulo comun para los informes de la ronda 3.
% Se incluye con \input{../comun/preambulo} desde cada subcarpeta.
\usepackage[utf8]{inputenc}
\usepackage[T1]{fontenc}
\usepackage[spanish,es-nodecimaldot]{babel}
\usepackage{lmodern}
\usepackage[a4paper,margin=2.4cm]{geometry}
\usepackage{booktabs}
\usepackage{longtable}
\usepackage{array}
\usepackage{xcolor}
\usepackage{hyperref}
\usepackage{enumitem}
\usepackage{fancyhdr}
\usepackage{titlesec}
\usepackage{graphicx}
\usepackage{amssymb}
\usepackage{microtype}

\definecolor{sosTeal}{HTML}{0F766E}
\definecolor{sosGray}{HTML}{475569}

\hypersetup{
  colorlinks=true,
  linkcolor=sosTeal,
  urlcolor=sosTeal,
  citecolor=sosTeal,
  pdftitle={SOS Clinic - Ronda 3 de pruebas},
  pdfauthor={Equipo SOS Clinic}
}

\pagestyle{fancy}
\fancyhf{}
\lhead{\textcolor{sosGray}{SOS Clinic --- Ronda 3}}
\rhead{\textcolor{sosGray}{\leftmark}}
\cfoot{\thepage}
\setlength{\headheight}{15pt}
\renewcommand{\headrulewidth}{0.4pt}

\titleformat{\section}{\Large\bfseries\color{sosTeal}}{\thesection}{0.8em}{}
\titleformat{\subsection}{\large\bfseries\color{sosGray}}{\thesubsection}{0.6em}{}

\setlist[itemize]{leftmargin=1.4em,itemsep=2pt,topsep=4pt}
\setlist[enumerate]{leftmargin=1.6em,itemsep=2pt,topsep=4pt}

\newcommand{\pantalla}[1]{\texttt{\detokenize{#1}}}
\newcommand{\ok}{\textcolor{sosTeal}{\textbf{PASS}}}
\newcommand{\fallo}{\textcolor{red!70!black}{\textbf{FAIL}}}
\newcommand{\aviso}{\textcolor{orange!80!black}{\textbf{WARN}}}
\newcommand{\omitido}{\textcolor{sosGray}{\textbf{SKIP}}}
 desde cada subcarpeta.
\usepackage[utf8]{inputenc}
\usepackage[T1]{fontenc}
\usepackage[spanish,es-nodecimaldot]{babel}
\usepackage{lmodern}
\usepackage[a4paper,margin=2.4cm]{geometry}
\usepackage{booktabs}
\usepackage{longtable}
\usepackage{array}
\usepackage{xcolor}
\usepackage{hyperref}
\usepackage{enumitem}
\usepackage{fancyhdr}
\usepackage{titlesec}
\usepackage{graphicx}
\usepackage{amssymb}
\usepackage{microtype}

\definecolor{sosTeal}{HTML}{0F766E}
\definecolor{sosGray}{HTML}{475569}

\hypersetup{
  colorlinks=true,
  linkcolor=sosTeal,
  urlcolor=sosTeal,
  citecolor=sosTeal,
  pdftitle={SOS Clinic - Ronda 3 de pruebas},
  pdfauthor={Equipo SOS Clinic}
}

\pagestyle{fancy}
\fancyhf{}
\lhead{\textcolor{sosGray}{SOS Clinic --- Ronda 3}}
\rhead{\textcolor{sosGray}{\leftmark}}
\cfoot{\thepage}
\setlength{\headheight}{15pt}
\renewcommand{\headrulewidth}{0.4pt}

\titleformat{\section}{\Large\bfseries\color{sosTeal}}{\thesection}{0.8em}{}
\titleformat{\subsection}{\large\bfseries\color{sosGray}}{\thesubsection}{0.6em}{}

\setlist[itemize]{leftmargin=1.4em,itemsep=2pt,topsep=4pt}
\setlist[enumerate]{leftmargin=1.6em,itemsep=2pt,topsep=4pt}

\newcommand{\pantalla}[1]{\texttt{\detokenize{#1}}}
\newcommand{\ok}{\textcolor{sosTeal}{\textbf{PASS}}}
\newcommand{\fallo}{\textcolor{red!70!black}{\textbf{FAIL}}}
\newcommand{\aviso}{\textcolor{orange!80!black}{\textbf{WARN}}}
\newcommand{\omitido}{\textcolor{sosGray}{\textbf{SKIP}}}
 desde cada subcarpeta.
\usepackage[utf8]{inputenc}
\usepackage[T1]{fontenc}
\usepackage[spanish,es-nodecimaldot]{babel}
\usepackage{lmodern}
\usepackage[a4paper,margin=2.4cm]{geometry}
\usepackage{booktabs}
\usepackage{longtable}
\usepackage{array}
\usepackage{xcolor}
\usepackage{hyperref}
\usepackage{enumitem}
\usepackage{fancyhdr}
\usepackage{titlesec}
\usepackage{graphicx}
\usepackage{amssymb}
\usepackage{microtype}

\definecolor{sosTeal}{HTML}{0F766E}
\definecolor{sosGray}{HTML}{475569}

\hypersetup{
  colorlinks=true,
  linkcolor=sosTeal,
  urlcolor=sosTeal,
  citecolor=sosTeal,
  pdftitle={SOS Clinic - Ronda 3 de pruebas},
  pdfauthor={Equipo SOS Clinic}
}

\pagestyle{fancy}
\fancyhf{}
\lhead{\textcolor{sosGray}{SOS Clinic --- Ronda 3}}
\rhead{\textcolor{sosGray}{\leftmark}}
\cfoot{\thepage}
\setlength{\headheight}{15pt}
\renewcommand{\headrulewidth}{0.4pt}

\titleformat{\section}{\Large\bfseries\color{sosTeal}}{\thesection}{0.8em}{}
\titleformat{\subsection}{\large\bfseries\color{sosGray}}{\thesubsection}{0.6em}{}

\setlist[itemize]{leftmargin=1.4em,itemsep=2pt,topsep=4pt}
\setlist[enumerate]{leftmargin=1.6em,itemsep=2pt,topsep=4pt}

\newcommand{\pantalla}[1]{\texttt{\detokenize{#1}}}
\newcommand{\ok}{\textcolor{sosTeal}{\textbf{PASS}}}
\newcommand{\fallo}{\textcolor{red!70!black}{\textbf{FAIL}}}
\newcommand{\aviso}{\textcolor{orange!80!black}{\textbf{WARN}}}
\newcommand{\omitido}{\textcolor{sosGray}{\textbf{SKIP}}}


\title{\textbf{Pentest local}\\[0.3em]
\large Metodología de las pruebas de seguridad --- Ronda 3}
\author{Equipo SOS Clinic}
\date{22 de octubre de 2026}

\begin{document}
\maketitle

\section{Objetivo y alcance}

Re-verificar los 25 hallazgos de la ronda 2 y comprobar la postura de seguridad
actual del sistema: autenticación, autorización por rol, JWT, CORS, cabeceras,
entrada de datos, IDOR, límites de intentos y XSS almacenado.

Todo el tráfico se dirigió a \texttt{localhost} (gateway 8080, servicios
8081--8084 y 8088, frontend 5500). No se atacó infraestructura externa. Las
pruebas se ejecutaron sobre el seed canónico con cuentas reales de rol y
cuentas temporales para los casos de fuerza bruta.

\section{Entorno}

\begin{itemize}
  \item Stack completo levantado con \pantalla{04. levantar-sistema.bat}.
  \item PostgreSQL con el seed de 9 cuentas, 24 pacientes y 30 episodios.
  \item Ocho roles disponibles: administrador, seguridad, coordinación,
  médico, triage, recepción, lista de espera y auditor.
  \item Navegador Chromium controlado para la prueba de XSS.
\end{itemize}

\section{Diseño de las pruebas}

\subsection{Autenticación y JWT}

\begin{itemize}
  \item Token firmado con un secreto distinto, token con \texttt{alg=none} y
  token real con el payload modificado (perfil elevado a administrador).
  \item Peticiones sin token al gateway y a los puertos directos de cada
  servicio, para comprobar que la validación no depende solo del gateway.
\end{itemize}

\subsection{Matriz de autorización}

Se probaron 11 endpoints con los 8 roles (88 combinaciones), verificando el
código esperado en cada caso: lecturas clínicas, escrituras, administración,
auditoría, aprobación de cuentas, analítica y agente IA.

\subsection{Configuración}

\begin{itemize}
  \item Preflight CORS con origen permitido y con \texttt{evil.local}.
  \item Cabeceras de seguridad del gateway y del frontend.
  \item CSP, traversal en el servidor estático y exposición de Actuator.
\end{itemize}

\subsection{Datos y lógica}

\begin{itemize}
  \item Mass assignment en pacientes y signos vitales (campos de servidor
  enviados por el cliente).
  \item Inyección SQL en búsqueda por cédula y cédula exacta.
  \item JSON malformado, referencias inexistentes y episodios inexistentes.
  \item IDOR horizontal: acceso a episodios, pacientes e historias de otros
  usuarios con un rol clínico.
\end{itemize}

\subsection{Límites de intentos y 2FA}

\begin{itemize}
  \item Challenge 2FA válido, reutilización y propósito incorrecto.
  \item Cinco códigos 2FA erróneos y bloqueo posterior.
  \item Cinco contraseñas erróneas y bloqueo posterior.
  \item Enumeración de usuarios en login y recuperación de clave.
\end{itemize}

\subsection{XSS almacenado}

Se creó una cuenta temporal sin 2FA y un paciente con el nombre
\texttt{<img src=x onerror=window.\_\_xss=1>}. Un navegador real abrió
\texttt{pacientes.html}, \texttt{lista-espera.html} y
\texttt{seguimiento.html?idRegistro=37}; se comprobó la variable
\texttt{window.\_\_xss}, la presencia de \texttt{<img>} inyectadas, el texto
visible y los diálogos.

\section{Procedimiento y limpieza}

\begin{enumerate}
  \item Login de las 8 cuentas del seed completando el 2FA con un código
  temporal insertado en la base (no se leen correos).
  \item Ejecución de las comprobaciones con \pantalla{pentest_local.py}.
  \item Creación de usuarios temporales únicos por corrida para evitar el
  bloqueo en memoria del limitador.
  \item Limpieza de usuarios, códigos, tokens, tomas y pacientes de prueba al
  finalizar, incluso si una sección falla.
  \item Prueba de XSS con \pantalla{pentest_xss.py}, navegador y limpieza
  posterior.
\end{enumerate}

El limitador de login/2FA es en memoria. Repetir el pentest en menos de 10
minutos devuelve 429; en local se reinició \texttt{ms-seguridad} para limpiar la
ventana antes del pase final.

\section{Criterios y salvaguardas}

\begin{itemize}
  \item \ok{} si el comportamiento coincide con la matriz esperada o con la
  remediación de la ronda 2.
  \item \fallo{} si hay una desviación de seguridad o un error no controlado.
  \item \aviso{} para residuales aceptados o riesgos de mínimo privilegio.
  \item No se imprimen credenciales, tokens ni contenido de \texttt{.env}.
  \item Los correos de prueba llegan a la bandeja
  \texttt{pruebasosclinic@gmail.com}; no se usan datos reales.
  \item No se ejecutan ataques de denegación de servicio ni fuerza bruta masiva.
\end{itemize}

\section{Limitaciones}

\begin{itemize}
  \item El IDOR se cierra por rol, no por propiedad del episodio; es un
  residual aceptado y se documenta como tal.
  \item No se probó revocación de JWT (el sistema no la tiene).
  \item No se evaluaron ataques de red, contenedores ni cadena de suministro
  más allá de Trivy.
  \item La prueba de XSS cubre tres páginas; no es un barrido exhaustivo de
  todas las vistas.
\end{itemize}

\section{Reproducción}

\begin{verbatim}
python seguridad/pruebas/pentest_local.py
python seguridad/pruebas/pentest_xss.py prepare
# abrir pacientes/lista de espera con la cuenta temporal
python seguridad/pruebas/pentest_xss.py cleanup
\end{verbatim}

\end{document}
